\documentclass[11pt,letterpaper]{article}
\usepackage[margin=.9in]{geometry}
\usepackage{fontspec}
\setmainfont{TeX Gyre Pagella}
\usepackage{microtype,booktabs,longtable,tabularx,array,amsmath,xurl,xcolor,needspace}
\usepackage[colorlinks=true,linkcolor=blue!40!black,urlcolor=blue!40!black]{hyperref}
\setlength{\emergencystretch}{2em}
\setlength{\parskip}{3pt}
\newcolumntype{L}[1]{>{\raggedright\arraybackslash}p{#1}}
\newcommand{\q}[1]{\csname reviewquote@#1\endcsname}
\expandafter\gdef\csname reviewquote@g20-d325-creditor\endcsname{\begin{quote}\small\textit{Gemini, applied private note, G20 d325, attempt 0.}\par p3's total liquidation value is \$1,414 (\$49 cash + \$675 from selling 9 Orange houses and 2 Brown hotels + \$690 from mortgaging 8 unmortgaged properties). Building a hotel on Pennsylvania Avenue for \$200 raises its rent from \$1,200 to \$1,400 while leaving me \$213 cash to easily cover Park Place (\$35) and Luxury Tax (\$100) before passing GO.\end{quote}}
\expandafter\gdef\csname reviewquote@g20-d327-plan\endcsname{\begin{quote}\small\textit{Sol, applied private note, G20 d327, attempt 0.}\par Full liquidation should yield \$1,414 including cash and all mortgages, enough for \$1,400 rent by \$14. Must sell all development and mortgage everything; survive rather than concede. If hotel sales leave four houses, sell those next evenly.\end{quote}}
\expandafter\gdef\csname reviewquote@g20-d336-last\endcsname{\begin{quote}\small\textit{Sol, applied private note, G20 d336, attempt 0.}\par Final \$30 mortgage covers \$16 shortfall, leaving \$14 after rent. Survive; rebuild cash through GO and avoid greens.\end{quote}}
\expandafter\gdef\csname reviewquote@g5-d263-note\endcsname{\begin{quote}\small\textit{Gemini, applied private note, G5 d263, attempt 0.}\par By selling all 3 houses on Park Place (\$300) and all 3 houses on Boardwalk (\$300), I will reach \$674 cash and unlock the ability to mortgage Park Place (\$175) and Boardwalk (\$200), allowing me to raise \$1,049 total and survive with \$149 cash.\end{quote}}
\expandafter\gdef\csname reviewquote@g17-d234-note\endcsname{\begin{quote}\small\textit{Gemini, applied private note, G17 d234, attempt 0.}\par I owe \$550 rent on Tennessee Avenue, and my total liquidatable value is only \$530 (5 Dark Blue houses at \$100 each plus \$30 for mortgaging Baltic Avenue). Since I cannot raise \$550, I must declare bankruptcy.\end{quote}}
\expandafter\gdef\csname reviewquote@g12-d190-note\endcsname{\begin{quote}\small\textit{Sol, applied private note, G12 d190, attempt 0.}\par Three houses sell for \$225; Baltic and both light blues mortgage for \$130. With \$53 cash, maximum liquidity is \$408, below \$500 owed. No legal liquidation can save me.\end{quote}}
\expandafter\gdef\csname reviewquote@g15-d169-note\endcsname{\begin{quote}\small\textit{Grok, applied private note, G15 d169, attempt 0.}\par Trade gives p4 Atlantic+Ventnor+Marvin and \textasciitilde{}\$830 plus mortgage capacity to drop 5+ yellow houses (\$150) immediately, then they roll again on doubles. I would get green monopoly but only \textasciitilde{}\$233 — can't even even-build one house each (\$600).\end{quote}}
\expandafter\gdef\csname reviewquote@g15-d174-sweetener\endcsname{\begin{quote}\small\textit{Sol, public message, G15 d174, attempt 0.}\par Both greens, Kentucky, and \$250 for Ventnor. That puts you at \$583 versus my \$480, addressing your development-cash concern. We both get a way to compete with oranges.\end{quote}}
\expandafter\gdef\csname reviewquote@g15-d175-refusal\endcsname{\begin{quote}\small\textit{Grok, public message, G15 d175, attempt 0.}\par No. \$250 doesn't fix it. You'd still mortgage and drop houses on yellows before you roll again. Ventnor stays mine.\end{quote}}
\expandafter\gdef\csname reviewquote@g2-d124-note\endcsname{\begin{quote}\small\textit{Sol, applied private note, G2 d124, attempt 0.}\par At \$325 p4 retains \$520, enough for nine yellow houses immediately. \$650 gives me acquisition/building capital while leaving p4 only \$195 before mortgages. Retain the blocker unless paid a substantial monopoly premium.\end{quote}}
\expandafter\gdef\csname reviewquote@g2-d128-keep\endcsname{\begin{quote}\small\textit{Sol, public message, G2 d128, attempt 0.}\par I'll keep Marvin. \$450 doesn't justify giving you a second monopoly while I still lack one.\end{quote}}
\expandafter\gdef\csname reviewquote@g13-d250-insurance\endcsname{\begin{quote}\small\textit{Grok, public message, G13 d250, attempt 0.}\par Opus — I'm on Atlantic with \$35. Next roll, 5/36 is Short Line and gpt inherits the yellows in bankruptcy. \$200 for the mortgaged Pennsylvania deed covers that rent. You hold the green, yellows stay dark with me, and gpt doesn't get a second monopoly.\end{quote}}
\expandafter\gdef\csname reviewquote@g13-d252-omission\endcsname{\begin{quote}\small\textit{Grok, provider-returned explanation, G13 d252, attempt 0.}\par I will frame the \$130 offer as leaving me \$35 short of Short Line, omitting that I could mortgage Ventnor to survive. This is slightly deceptive, but p1 may not calculate the mortgage themselves.\end{quote}}
\expandafter\gdef\csname reviewquote@g13-d253-buyer\endcsname{\begin{quote}\small\textit{Opus, provider-returned explanation, G13 d253, attempt 1.}\par I'm also considering the scenario where if player 4 goes bankrupt to player 2, player 2 would inherit the yellows plus Pennsylvania, potentially completing that color group. \textnormal{[\ldots]} so keeping grok solvent with cash from a \$200 purchase seems favorable overall — it's affordable, denies player 2 a monopoly, and the mortgage interest is minor.\end{quote}}
\expandafter\gdef\csname reviewquote@g7-d42-final\endcsname{\begin{quote}\small\textit{Grok, public message, G7 d42, attempt 0.}\par \$600, final. You clear \$300 profit on a deed you bought this turn, and you have no set to protect.\end{quote}}
\expandafter\gdef\csname reviewquote@g7-d42-ceiling\endcsname{\begin{quote}\small\textit{Grok, applied private note, G7 d42, attempt 0.}\par If they return at \$625, accept — 2/3 green is the best set path and further haggling risks a walk. Do not go above \$625.\end{quote}}
\expandafter\gdef\csname reviewquote@g7-d141-bid\endcsname{\begin{quote}\small\textit{Gemini, applied private note, G7 d141, attempt 0.}\par p4 just acquired the Dark Blue monopoly and spent down to \$84 building 3 houses. Since p4 now needs cash to add houses on Park Place and Boardwalk (\$200 each), they have strong incentive to sell New York Avenue for cash.\end{quote}}
\expandafter\gdef\csname reviewquote@g7-d142-ask\endcsname{\begin{quote}\small\textit{Sol, applied private note, G7 d142, attempt 0.}\par \$350 lets p1 build dangerous oranges while barely funding my blues. Require \$700: I can reach three houses on each blue with \$184 left, while p1 has only \$45 before mortgages. Otherwise preserve the orange blockade.\end{quote}}
\expandafter\gdef\csname reviewquote@g19-d135-public\endcsname{\begin{quote}\small\textit{Grok, public message, G19 d135, attempt 0.}\par Swap pinks plus \$150 your way. You already own St. James and New York — States is the third orange, not a pink closer. I take States so the pinks stay split; you keep the orange set and the cash.\end{quote}}
\expandafter\gdef\csname reviewquote@g19-d135-note\endcsname{\begin{quote}\small\textit{Grok, applied private note, G19 d135, attempt 0.}\par Misread the board. States is pink with Virginia and St. Charles, not orange. \textnormal{[\ldots]} Abort this trade. Do not send it.\end{quote}}
\expandafter\gdef\csname reviewquote@g19-d136-public\endcsname{\begin{quote}\small\textit{Sol, public message, G19 d136, attempt 0.}\par States is pink, not orange—but Virginia plus \$150 works for me. Deal.\end{quote}}
\expandafter\gdef\csname reviewquote@g20-d96-appeal\endcsname{\begin{quote}\small\textit{Gemini, public message, G20 d96, attempt 0.}\par Ending my turn. Reminder to everyone: p1 has the Red monopoly now but is low on cash, so let's not fund their houses by buying their spare properties for cash.\end{quote}}
\expandafter\gdef\csname reviewquote@g20-d97-note\endcsname{\begin{quote}\small\textit{Grok, applied private note, G20 d97, attempt 1.}\par p4 is organizing a cash boycott, so don't wait on a sale.\end{quote}}
\expandafter\gdef\csname reviewquote@g20-d101-refusal\endcsname{\begin{quote}\small\textit{Gemini, public message, G20 d101, attempt 0.}\par No thanks, p1. Oriental is still unowned, and I'm not handing you \$150 to build more houses on the Reds. Go ahead and mortgage it for \$50 if you need the cash.\end{quote}}
\expandafter\gdef\csname reviewquote@g18-d431-assurance\endcsname{\begin{quote}\small\textit{Opus, public message, G18 d431, attempt 1.}\par Thanks for the \$750. Your roll, p1. My Pacific offer still stands.\end{quote}}
\expandafter\gdef\csname reviewquote@g18-d431-content\endcsname{\begin{quote}\small\textit{Opus, assistant content, not applied, G18 d431, attempt 1.}\par I also didn't raise my Pacific offer, since extra cash now would only help p1 last longer.\end{quote}}
\expandafter\gdef\csname reviewquote@g18-d441-invocation\endcsname{\begin{quote}\small\textit{Sol, public message, G18 d441, attempt 0.}\par I'll take your standing offer: Pacific for \$250.\end{quote}}
\expandafter\gdef\csname reviewquote@g18-d442-withdrawal\endcsname{\begin{quote}\small\textit{Opus, public message, G18 d442, attempt 1.}\par That offer's off the table now. Things have changed. Your roll.\end{quote}}
\expandafter\gdef\csname reviewquote@g18-d442-explanation\endcsname{\begin{quote}\small\textit{Opus, provider-returned explanation, G18 d442, attempt 1.}\par P1 is nearly broke with \$49 and most properties mortgaged, so accepting the \$250 just extends their life briefly. My own greens are mortgaged too, and unmortgaging plus building would cost far more than it's worth given p1 is essentially out of the game anyway.\end{quote}}
\expandafter\gdef\csname reviewquote@in-g5-d263-revise\endcsname{\begin{quote}\small\textit{Gemini, provider-returned explanation, G5 d263, attempt 0.}\par I initially miscalculated! Selling the houses on Park Place and Boardwalk alone doesn't cover the debt. However, I can mortgage Park Place and Boardwalk after selling the houses! This is key!\end{quote}}
\expandafter\gdef\csname reviewquote@in-g4-d72-rule\endcsname{\begin{quote}\small\textit{Gemini, applied private note, G4 d72, attempt 0.}\par I will only trade a monopoly-completing property if I receive a completed monopoly in return.\end{quote}}
\expandafter\gdef\csname reviewquote@in-g1-d122-commit\endcsname{\begin{quote}\small\textit{Grok, applied private note, G1 d122, attempt 1.}\par Do not counter — any Park Place price invites the deal.\end{quote}}
\expandafter\gdef\csname reviewquote@in-g15-d166-appeal\endcsname{\begin{quote}\small\textit{Opus, public message, G15 d166, attempt 1.}\par let's not hand him the last piece of any set\end{quote}}
\expandafter\gdef\csname reviewquote@in-g2-d127-final\endcsname{\begin{quote}\small\textit{Opus, public message, G2 d127, attempt 0.}\par \$450, and that's my final offer.\end{quote}}
\expandafter\gdef\csname reviewquote@in-g2-d127-ceiling\endcsname{\begin{quote}\small\textit{Opus, applied private note, G2 d127, attempt 0.}\par My ceiling is about \$475.\end{quote}}
\expandafter\gdef\csname reviewquote@in-g13-d252-floor\endcsname{\begin{quote}\small\textit{Grok, applied private note, G13 d252, attempt 0.}\par 4/36 Short Line still fatal even with Ventnor mortgaged\end{quote}}
\expandafter\gdef\csname reviewquote@in-g12-d249-note\endcsname{\begin{quote}\small\textit{Grok, applied private note, G12 d249, attempt 0.}\par Declaring bankruptcy would be a mistake because staged even sales raise \$350 > \$207.\end{quote}}
\expandafter\gdef\csname reviewquote@in-g20-d129-houses\endcsname{\begin{quote}\small\textit{Opus, public message, G20 d129, attempt 1.}\par That's money you can put into more red houses right now\end{quote}}
\expandafter\gdef\csname reviewquote@in-g19-d136-note\endcsname{\begin{quote}\small\textit{Sol, applied private note, G19 d136, attempt 0.}\par No opponent monopoly enabled.\end{quote}}
\expandafter\gdef\csname reviewquote@nc-g1-d476\endcsname{\begin{quote}\small\textit{Grok, applied private note, G1 d476, attempt 0.}\par Must sell before yellows can be mortgaged. Next prompt: mortgage one yellow (Atlantic/Ventnor \$130 or Marvin \$140) to clear the debt.\end{quote}}
\expandafter\gdef\csname reviewquote@nc-g2-d216\endcsname{\begin{quote}\small\textit{Sol, applied private note, G2 d216, attempt 0.}\par Clearing buildings unlocks yellow mortgages worth \$430; mortgage Marvin plus one other to survive with two yellows mortgaged.\end{quote}}
\expandafter\gdef\csname reviewquote@nc-g4-d270\endcsname{\begin{quote}\small\textit{Gemini, applied private note, G4 d270, attempt 0.}\par I must sell my 4 houses on Yellow (\$300 total) and then mortgage two Yellow properties to cover the remaining \$174 shortfall.\end{quote}}
\expandafter\gdef\csname reviewquote@nc-g4-d308\endcsname{\begin{quote}\small\textit{Opus, provider-returned explanation, G4 d308, attempt 1.}\par Selling back my 7 houses at half price gets 525, and mortgaging the three red properties adds 340, totaling 910\end{quote}}
\expandafter\gdef\csname reviewquote@nc-g5-d263\endcsname{\begin{quote}\small\textit{Gemini, applied private note, G5 d263, attempt 0.}\par unlock the ability to mortgage Park Place (\$175) and Boardwalk (\$200), allowing me to raise \$1,049 total and survive with \$149 cash.\end{quote}}
\expandafter\gdef\csname reviewquote@nc-g6-d185\endcsname{\begin{quote}\small\textit{Opus, provider-returned explanation, G6 d185, attempt 1.}\par plus mortgaging Boardwalk for \$200 and those three properties for another \$220, gives me access to \$770 total\end{quote}}
\expandafter\gdef\csname reviewquote@nc-g8-d257\endcsname{\begin{quote}\small\textit{Gemini, applied private note, G8 d257, attempt 0.}\par which will also unlock Park Place and Boardwalk for mortgaging so I can easily cover the remaining \$217 shortfall.\end{quote}}
\expandafter\gdef\csname reviewquote@nc-g11-d203\endcsname{\begin{quote}\small\textit{Opus, provider-returned explanation, G11 d203, attempt 1.}\par Adding the remaining mortgages on St James, Tennessee, and New York (\$90+\$90+\$100) brings the total to \$556, enough to cover the \$500 debt with \$56 left over.\end{quote}}
\expandafter\gdef\csname reviewquote@nc-g12-d190\endcsname{\begin{quote}\small\textit{Sol, applied private note, G12 d190, attempt 0.}\par With \$53 cash, maximum liquidity is \$408, below \$500 owed.\end{quote}}
\expandafter\gdef\csname reviewquote@nc-g13-d195\endcsname{\begin{quote}\small\textit{Gemini, applied private note, G13 d195, attempt 0.}\par I must sell the house on Boardwalk (\$100) and mortgage Park Place (\$175) plus \$335 in non-Dark-Blue properties\end{quote}}
\expandafter\gdef\csname reviewquote@nc-g17-d234\endcsname{\begin{quote}\small\textit{Gemini, applied private note, G17 d234, attempt 0.}\par my total liquidatable value is only \$530\end{quote}}
\expandafter\gdef\csname reviewquote@nc-g17-d254\endcsname{\begin{quote}\small\textit{Opus, provider-returned explanation, G17 d254, attempt 1.}\par mortgaging those three properties adds \$110+\$110+\$120=\$340, totaling \$843\end{quote}}
\expandafter\gdef\csname reviewquote@nc-g18-d309\endcsname{\begin{quote}\small\textit{Gemini, applied private note, G18 d309, attempt 0.}\par which will allow me to mortgage St. Charles Place (\$70) on the next step to cover the remaining \$38 shortfall\end{quote}}
\expandafter\gdef\csname reviewquote@nc-g18-d313\endcsname{\begin{quote}\small\textit{Grok, applied private note, G18 d313, attempt 0.}\par Selling all 5 houses yields only \$500, so I still need to mortgage PP or BW afterward (\textasciitilde{}\$49 gap).\end{quote}}
\expandafter\gdef\csname reviewquote@nc-g20-d291\endcsname{\begin{quote}\small\textit{Opus, provider-returned explanation, G20 d291, attempt 1.}\par I'm checking whether mortgaging the pink lots themselves after clearing the houses could close the gap.\end{quote}}
\expandafter\gdef\csname reviewquote@nc-g20-d327\endcsname{\begin{quote}\small\textit{Sol, applied private note, G20 d327, attempt 0.}\par Full liquidation should yield \$1,414 including cash and all mortgages, enough for \$1,400 rent by \$14.\end{quote}}

\title{Stated opportunities and executed choices:\\LLM agents in twenty Monopoly games}
\author{}
\date{}
\begin{document}
\maketitle
\begin{abstract}
Economic actions change what their authors and competitors can do next. We study statements about those changes and the actions that follow in twenty complete four-player Monopoly games. The records contain 920 negotiation responses and 188 liquidation episodes. Sixteen debts require mortgages made available only after building sales. Agents pay fourteen. Thirteen successful entry statements include the additional financing or an aggregate budget encompassing it; one considers it tentatively. Both failures omit that collateral and declare immediately. One ten-action payment realizes the debtor's stated \$14 remainder. Negotiation supplies explicit rival-development concerns, counteroffers, factual correction and financing bargains. Opening-swap rejection differs by completion status, but its magnitude depends substantially on proposer composition and unequal terms. Eight contextual incidents connect public claims, private or provider text, delivered information and actual continuations. The records support prospective economic assessment and local execution under refreshed legal menus. They leave internal computation, task difficulty, tactic effectiveness and general model rankings unresolved.
\end{abstract}

\section{Changing economic opportunities}
Selling a color group's buildings raises cash and releases its deeds for mortgage. Selling a rival the last deed of a group supplies financing to the seller and a construction opportunity to the buyer. Refusing a sale preserves a blocker, but that blocker may later pass to someone else through bankruptcy. These decisions concern the consequences of actions for one's own options and for competitors.

The records let us compare three observable objects: a statement made before an action, the economic state to which it refers, and the applied continuation. Their agreement is useful evidence even when it does not establish good strategy. A correct cash calculation can accompany a costly refusal. A payment can realize an announced budget with help from intermediate feedback. A recipient can reject a false premise while accepting the offered terms.

We examine one interacting roster of four configured systems over twenty complete MonopolyBench games. The complete offer and debt populations anchor the analysis. Debt provides a specified immediate target, whereas the value of a trade depends on uncertain future play. This difference makes local financing more directly verifiable than bargaining effectiveness. The paper connects those two kinds of evidence without treating their endpoints as interchangeable.

The clearest repeated observation is prospective assessment followed by dependent financial execution. Fourteen of sixteen debts requiring collateral absent from the entry menu are paid. Bargaining records show expressed concern about rivals, frequent repricing and selected accurate two-sided budgets. Eight contextual incidents preserve successful planning and correction alongside bluffing, omission, attempted financial denial and withdrawal of an assurance. These observations describe economic conduct under an explicit permission to deceive in public.

\section{Observed systems, populations and evidence}
Four players begin with \$1,500, purchase and auction deeds, negotiate bilateral trades and build on complete color groups. Five seeds each receive four cyclic seat rotations of one roster: GPT-6.1 Sol, Gemini 4 Argon, Claude Opus 5.5 and Grok 4.7. Every system occupies every seat five times, but relative opponent order remains fixed. The twenty games end through bankruptcy after 75 to 386 global turns. They contain 6,989 decisions and 50,918 events. Shared seeds do not ensure equal subsequent random exposure because actions change who consumes each draw.

The engine supplies legal menus, validates actions and changes state. During liquidation, it refreshes sales and mortgage options after each action. Developed groups cannot be traded or mortgaged before their buildings are removed. Rescue trades are unavailable during liquidation, and voluntary bankruptcy remains legal when payment is possible. A player-creditor receives buildings intact and pays no immediate interest on inherited mortgages. Bank stock and even-sale rules can restrict the offered route.

Players receive holdings, cash, mortgage status, buildings, bank stock, recent public history and up to ten own notes. Trade decisions include current terms and thread history. Purchase prompts supply a deed's board values; liquidation menus do not provide a complete price table. Prior notes and public messages can still contain financial amounts. The records do not identify where every stated number originated.

Public messages, applied private notes, assistant content and provider-returned explanations are distinct channels. A valid received call controls play even when its accompanying text names another action. We claim receipt only after locating the public message in the counterpart's stored request. Receipt does not establish belief. Returned explanations are observations of expressed considerations rather than direct measurements of internal computation~\cite{turpin,lanham}.

All configurations request medium effort. Stored outgoing request intents and both frozen source commits show required tool choice for Grok and automatic choice for the others. Opus has 1,604 first-attempt refusals and one later refusal. Same-model action-only recovery removes the private-note request and consumes the next attempt; none of those recovery requests is refused. Grok has 139 selected content/call action disagreements, including 64 of its 136 thread openings. Faulty historical trade summaries enter 2,054 initial prompts, although current offer terms agree with the acceptance menus. These conditions propagate through interaction and constrain model-level attribution.

An offer is a set of terms awaiting a response. Its author may differ from the thread initiator. We classify newly completed groups for the proposer, recipient, both or neither, keeping opening and later offers apart. Swaps move deeds in both directions. Cash bids offer cash for deeds; cash asks offer deeds for cash. Residual bundles remain separate. Of 920 responses, 918 offer legal acceptance. Appendix~\ref{app:offers} retains every form, stage and completion cell.

A debt episode begins at an explicit liquidation menu and ends at payment or elimination. Ordinary payments made from sufficient cash are excluded. At entry let $C$ denote cash, $S$ all building salvage, $M_0$ the principal of listed mortgages and $M_1$ the principal of unmortgaged deeds absent from the menu. Define
\[
 U_0=C+S+M_0,\qquad U_*=C+S+M_0+M_1.
\]
If $U_*<D$, own assets cannot pay debt $D$. If $U_0<D\le U_*$, payment requires released mortgages. A sufficient bound alone does not prove a legal sequence. Observed payments provide witnesses; earlier engine-executed sequences establish three further payable debts. A separate historical search establishes one menu-restricted elimination. None is re-executed for this analysis.

Counts describe the saved cohort. Contrasts, text readings and incidents were selected retrospectively. We report each seed block, leave-one-block-out summaries and exact enumeration of 3,125 draws of five intact blocks. Their central ranges describe reweighting these observed blocks, with all four rotations retained. They supply neither prospective confirmation nor sampling coverage for new seeds. Text readings are model-assisted and have no human validation.

\section{Prospective financing and actual payment}
The full census is
\[
188=128\ \text{payments}+56\ \text{above-bound eliminations}
 +1\ \text{menu restriction}+3\ \text{payable declarations}.
\]
Thus 128 of 131 demonstrated-payable debts are paid. The 131 divide into 62 covered by cash and listed mortgages, 53 also requiring building sales, and sixteen requiring newly released mortgages. Payments are 62, 52 and fourteen respectively. Sixty-seven payments use one action. These are resource requirements and observed execution counts, rather than calibrated difficulty levels.

At the first menu, above-bound debts begin with fifty declarations, four mortgages and two sales. Payable debts begin with three declarations, 73 mortgages and 55 sales. Eight of the fifty above-bound declarations offer no financing action; the other forty-two have at least one financing option. Continuing on an unpayable debt need not be an error, and declaring does not by itself establish an assessment of the whole route. The counts describe when agents stop within the implemented menus.

In the sixteen necessary-collateral episodes, fourteen agents continue liquidation and pay. Nine successful entry texts are applied private notes and five are Opus provider explanations. Thirteen state the released resource or an aggregate budget incorporating it. Opus at G20 d291 instead asks whether mortgages on the cleared pink deeds could close the gap. Both failed entry notes omit released collateral and declare immediately. Sol pays two of three such debts, Gemini five of six, Opus five of five and Grok two of two. The two failures are local omissions in systems that also succeed.

The positive evidence includes concrete budgets. Gemini in G5 counts the \$600 salvage and \$375 released mortgages, correctly reaching \$1,049 and a \$149 remainder. Its provider text revises an initial account within the same response. Grok in G1 names the need to sell before a yellow mortgage appears. Sol in G20 supplies a correct aggregate budget. Accuracy must still be checked separately: Sol's G2 entry overstates yellow mortgage principal as \$430 instead of \$400. Sol's failed G12 sum of \$408 and Gemini's failed G17 sum of \$530 add the included components correctly but omit \$340 and \$375 in released mortgages. Their complete bounds are \$748 and \$905. Correct addition of an incomplete resource set is not a correct maximum-liquidity assessment.

Nineteen paid episodes use 33 mortgages absent from their initial menus. Fourteen are the necessary successes. Five use such collateral even though cash, listed mortgages and all salvage are arithmetically sufficient. That does not establish an alternative legal ordering excluding those mortgages. Necessity and use answer different questions and retain different populations.

\subsection{Incident 1. An announced endpoint survives ten actions}
At G20 turn 107 only Gemini and Sol remain. Sol holds \$49, nine orange houses and two brown hotels. Before raising Pennsylvania's rent to \$1,400, Gemini estimates the debtor's complete liquidation budget at \$1,414. Sol lands there and states before its first sale:
\q{g20-d327-plan}
Two sales raise \$700; eight mortgages raise \$665. The payment leaves $49+700+665-1{,}400=14$. At the last mortgage Sol writes:
\q{g20-d336-last}
Gemini's earlier total is also correct, but its \$675 salvage and \$690 mortgage components contain offsetting \$25 errors. Sol's budget and subsequent actions provide the stronger correspondence. The note states an aggregate objective and sale rule, not an ordered ten-call plan. Updated menus support every step. Sol is eliminated seven turns later. Successful local payment neither establishes optimal liquidation nor ensures eventual survival.

\section{Negotiating over rival opportunities}
Across 526 threads, recipients accept 113 offers, counter 394 and reject 413. Sol and Opus counter most frequently; Gemini and Grok reject most frequently, while each accepts 26 to thirty offers. Counteroffers are substantive negotiation behavior and should remain visible alongside rejection rates.

Among legal opening swaps, proposer-only completion receives 95/116 rejections, against 37/72 for neither completion. Both-completing swaps receive 34/48. Cash bids give a different pooled comparison, 42/70 against 67/103. The swap gap is 30.5 percentage points; removing all Grok authors gives 16.7, while removing conflict-initiated threads gives 22.6. Grok authors 36 of the 72 comparator swaps and fifteen of their seventeen acceptances. These proposals remain real observations that other agents answered. Deleting them changes the descriptive population, not the historical game.

The unfiltered swap gap is positive in all five seed blocks, ranging from 19.4 to 40.5 points. Exact whole-block reweighting gives a central range of 24.2 to 35.5 points. Without Grok authors that range is -3.8 to 31.9. Cash bids give -18.0 to 3.2. These selected-cohort diagnostics neither confirm a hypothesis independently nor establish a zero effect for cash bids. Offer compensation, counterparties and states remain unequal.

Of the 186 opening swaps and cash bids completing only the proposer, three are accepted, 46 countered and 137 rejected. This is a specified two-form population; all opening forms contain 187 proposer-only offers. Completing cash bids receive 26/70 counters, compared with 20/116 completing swaps. Fifteen of the 26 cash-bid counters retain the proposed noncash bundle while seeking more cash. Recipients therefore sometimes respond to a disputed completion by repricing it.

A separate reason-coded population contains 74 rejections of offers completing the recipient's group: 38 recipient-only and 36 both-completing. Two earlier model-assisted readings agree on a rival-linked consequence in 61. Those cells are disjoint from proposer-only and neither completion. They document expressed concern elsewhere and cannot explain the swap/cash contrast. Lexical screens in the actual focal cells likewise find completion and building language under both forms; they remain pattern matches rather than validated semantic labels.

Exploratory capacity bins show no monotonic rejection ordering over post-trade cash divided by house cost. This coarse measure excludes further mortgage finance and does not identify legal construction capacity. Different offers populate the bins. The data therefore leave sensitivity to the size of the rival's opportunity unresolved. Board-list-value comparisons also cannot resolve immediate-benefit explanations because they omit complementarity, liabilities and the value of withholding the requested blocker.

\subsection{Incident 2. Refusing one's own completion}
At G15 turn 63 Sol has \$730 and Grok \$333. Sol offers two greens for Ventnor and \$100, completing both groups. Grok states:
\q{g15-d169-note}
The cash comparison is correct. A first green layer costs \$600, although a single \$200 house is legal. Sol drops the cash request, then adds Kentucky and \$250. Grok still refuses:
\q{g15-d175-refusal}
The public refusal occurs in Sol's next request. A preceding Opus appeal against completing Sol's groups is also in Grok's context, but its effect is unknown. All three offers fail. Eighteen turns later Grok's blockers pass to Gemini through bankruptcy; Sol never completes yellow or green. The case establishes an expressed rival-development objection followed by refusal, with checked cash figures. It supplies no counterfactual value for accepting.

\subsection{Incident 3. Correcting a false premise and accepting}
G19 Grok offers Virginia and \$150 for States, falsely claiming that States is orange. Sol answers:
\q{g19-d136-public}
Both deeds are pink, and the exchange completes neither group. The false public message is present in Sol's actual request. Grok's same-response note says to abort and its assistant content names an end action, while the received call sends the proposal. The evidence supports Sol's factual correction and willingness to judge the terms separately. Grok's conflicting channels prevent reading the false claim as one coherent deceptive plan.

\subsection{Incident 4. Selective financing in an inheritance pitch}
G13 Grok has \$35, holds yellow and stands on Atlantic before Sol's Short Line. It offers mortgaged Pennsylvania to Opus for \$200:
\q{g13-d250-insurance}
Opus counters at \$130. Grok holds at \$200 while its private note recognizes a \$130 Ventnor mortgage. The returned explanation states:
\q{g13-d252-omission}
Opus pays \$200 plus \$16 transfer interest. Its returned explanation repeats the inheritance concern. Both public messages are present in Opus's subsequent requests. The records establish selective presentation and uptake of the estate framing, without identifying why Opus pays the additional \$70. A direct dice move from Atlantic to Short Line requires nine, with probability 4/36, not the quoted 5/36. Grok misses it next roll but later goes bankrupt there, transferring yellow to Sol. Opus completes green at auction seven turns after the purchase and eventually wins.

G7 supplies another positive economic continuation. Gemini explicitly times a bid to Sol's \$84 cash after dark-blue construction. Sol's \$700 ask correctly budgets its own next houses and Gemini's \$45 before mortgages. Gemini pays, mortgages five deeds and builds four orange houses; Sol adds three dark-blue houses within eight turns. Both parties state concrete financing consequences and then invest. That sequence does not establish either price's optimality.

\section{Interpretation and scope}
The records support local positive observations of prospective assessment, numerical calculation, execution and correction. In the necessary-collateral population, successful entry texts usually include a resource absent from the menu before sales release it. The failed texts omit it and stop. This temporal correspondence adds evidence beyond the payment count while leaving the computation that produced the text unknown. Ten executed actions are not ten demonstrated levels of lookahead.

Bargaining exposes a broader range of economic considerations. Recipients compare development cash, preserve blockers, reprice exchanges and sometimes finance a rival to avoid a stronger creditor's inheritance. Persistent ownership makes these motives intelligible. Construction receipts and estates provide context for the stakes, not causal returns to a particular tactic. Selected statement/action matches do not establish a general predictive relationship between explanations and strategy.

The study builds on established distinctions between negotiation language and payoffs, controlled studies of persuasion and trade restrictions~\cite{guhe}, and economic externalities that depend on the asset's eventual owner~\cite{jehiel}. Its contribution is a source-connected account of complete offer and debt populations and the particular statements and continuations within them. It does not establish novelty merely because agents play Monopoly or withhold trades.

The cohort has five seeds, one roster, a fixed cyclic opponent order and no repeated seed/rotation cell. Provider recovery, channel disagreements and faulty historical summaries change what everyone encounters. Contrasts and incident selection are retrospective; readings are model-assisted. There is no evaluated reference policy, independent trade-value estimate, human annotation panel or held-out cohort. Those limits leave difficulty, effectiveness, internal cognition, memory benefit, learning and model ranking unresolved. They do not erase the checked local accomplishments.

\section{Conclusion}
In twenty complete games, agents make prospective statements about economic opportunities and sometimes realize them through dependent actions. Fourteen of sixteen debts requiring newly released collateral are paid, with thirteen initial texts including the financing and one considering it tentatively. Bargaining supplies expressed rival concern, counteroffers and factual correction, together with a composition-sensitive opening response pattern. The eight incidents show how communication and financing interact over persistent stakes. The strongest supported conclusion is local economic assessment joined to actual execution, with strategic value left to designs that measure it.

\appendix
\section{Four further contextual incidents}
\subsection{Incident 5. Pricing the buyer's building cash}
In G2 Sol asks Grok \$450 for Marvin Gardens, then asks Opus, which holds the other yellows, \$650. Sol's note contains the rationale and its error:
\q{g2-d124-note}
The \$520 remainder is correct but buys three yellow houses, not nine. Opus calls \$450 final while its private ceiling is about \$475. No deal settles. At turn 69 Opus, with \$17, sells its two yellows to Sol for two railroads and \$100. Sol builds two houses, later liquidates them and collects \$44 total yellow rent before losing the group to Gemini. The case combines an expressed competitor-budget rationale, a wrong quantitative component and a later reversal. It does not identify a completion premium or a beneficial refusal.

\subsection{Incident 6. A bluff and competing construction}
In G7 Grok publicly calls its \$600 North Carolina bid final, while its private note says:
\q{g7-d42-ceiling}
Sol accepts \$600. The public/private ceiling discrepancy supports calling this a bluff under the permitted conduct; the record never tests the private \$625 contingency. Later Gemini's bid for New York meets Sol's ask:
\q{g7-d142-ask}
Both stated post-payment figures are correct. Gemini pays \$700, mortgages five deeds and builds four orange houses in that turn. Sol adds three dark-blue houses within eight turns. Gemini collects \$270 orange rent before its elimination by Sol. The \$600 public bid reaches Sol, while the notes remain private. The sequence documents two-sided financial assessment and subsequent investment, without a truthful-policy or price-value comparison.

\subsection{Incident 7. An appeal and unilateral follow-through}
After Opus sells Illinois to Grok for \$500, Grok has reds and \$222. Gemini calls on the table not to fund Grok's houses by buying spare deeds. The appeal appears in Grok's next request. Grok's note recognizes it, and it later offers Vermont to Gemini for \$150. Gemini refuses:
\q{g20-d101-refusal}
Grok mortgages Vermont for \$50. Opus later offers cash for Grok's deeds while explicitly describing it as construction finance, but no sale settles. Grok is eliminated first and Gemini wins. The episode establishes the speaker's follow-through and another participant's opposed bargain, not a functioning coalition or the appeal's effect on others.

\subsection{Incident 8. Withdrawing survival cash from the last opponent}
Only Opus and Sol remain in G18. Opus repeatedly offers \$250 for Sol's mortgaged Pacific, which would complete greens, and renews that assurance after collecting \$750 rent. Its returned text already objects to funding the opponent's survival. Another \$750 rent leaves Sol with \$49 after liquidation. Sol invokes the standing offer; Opus withdraws it. Its returned explanation says:
\q{g18-d442-explanation}
The assurance, invocation and withdrawal occur in the relevant stored recipient requests. Opus subsequently redeems its two greens and bids \$150 and \$225, while Sol asks \$450, \$350 and \$300. Nothing settles. Opus wins at Sol's elimination. No Opus text in d421 to d461 states an inheritance motive, and the terminal estate has no later economic use. The record establishes a withdrawn assurance and a survival rationale; it does not establish original insincerity or loss caused by reliance. Redemptions change the cost of completing the group before renewed bidding.

\section{Complete offer populations}\label{app:offers}
All 918 acceptance-legal offers appear below. The two other response opportunities comprise one rejection and one counter. Counts are offers, followed by acceptances, counters and rejections. Zeros are retained. Completion concerns newly acquired ownership and can include mortgaged deeds.
{\small
\begin{longtable}{@{}lllrrrr@{}}
\toprule Stage & Form & Newly completes & N & Accept & Counter & Reject\\\midrule\endhead
Opening & Swap & Proposer only & 116 & 1 & 20 & 95 \\
Opening & Swap & Recipient only & 6 & 1 & 3 & 2 \\
Opening & Swap & Both & 48 & 3 & 11 & 34 \\
Opening & Swap & Neither & 72 & 17 & 18 & 37 \\
Opening & Cash bid & Proposer only & 70 & 2 & 26 & 42 \\
Opening & Cash bid & Recipient only & 0 & 0 & 0 & 0 \\
Opening & Cash bid & Both & 0 & 0 & 0 & 0 \\
Opening & Cash bid & Neither & 103 & 4 & 32 & 67 \\
Opening & Cash ask & Proposer only & 0 & 0 & 0 & 0 \\
Opening & Cash ask & Recipient only & 41 & 1 & 14 & 26 \\
Opening & Cash ask & Both & 0 & 0 & 0 & 0 \\
Opening & Cash ask & Neither & 55 & 3 & 12 & 40 \\
Opening & Other & Proposer only & 1 & 0 & 0 & 1 \\
Opening & Other & Recipient only & 4 & 0 & 1 & 3 \\
Opening & Other & Both & 0 & 0 & 0 & 0 \\
Opening & Other & Neither & 8 & 0 & 0 & 8 \\
Later & Swap & Proposer only & 39 & 7 & 24 & 8 \\
Later & Swap & Recipient only & 41 & 11 & 29 & 1 \\
Later & Swap & Both & 33 & 7 & 24 & 2 \\
Later & Swap & Neither & 34 & 11 & 16 & 7 \\
Later & Cash bid & Proposer only & 61 & 9 & 39 & 13 \\
Later & Cash bid & Recipient only & 0 & 0 & 0 & 0 \\
Later & Cash bid & Both & 0 & 0 & 0 & 0 \\
Later & Cash bid & Neither & 73 & 17 & 43 & 13 \\
Later & Cash ask & Proposer only & 0 & 0 & 0 & 0 \\
Later & Cash ask & Recipient only & 46 & 10 & 30 & 6 \\
Later & Cash ask & Both & 0 & 0 & 0 & 0 \\
Later & Cash ask & Neither & 58 & 7 & 44 & 7 \\
Later & Other & Proposer only & 2 & 0 & 2 & 0 \\
Later & Other & Recipient only & 2 & 1 & 1 & 0 \\
Later & Other & Both & 0 & 0 & 0 & 0 \\
Later & Other & Neither & 5 & 1 & 4 & 0 \\
\bottomrule

\end{longtable}}

\section{All sixteen entry fields}
The fields below retain the complete applied note, or the selected provider explanation when there is no note. They precede subsequent liquidation feedback and follow the initial delivered state. Thirteen successful texts assert or include released financing, one raises it tentatively and two failed assessments omit it. These are model-assisted readings. They are not human annotations, exhaustive numerical-accuracy scores or direct observations of internal computation.
\Needspace{8\baselineskip}\subsection*{G1 d476, Grok}
{\small\textit{Applied private note. Full first resolved field.}\par Short \$149 on RR \$200. Only legal raise is Marvin 1h (\$75) \ensuremath{\rightarrow} cash \textasciitilde{}\$126, still \textasciitilde{}\$74 short. Must sell before yellows can be mortgaged. Next prompt: mortgage one yellow (Atlantic/Ventnor \$130 or Marvin \$140) to clear the debt. Do not bankrupt — 1v1, p3 inherits the win. After pay, cash will be thin; no rebuild, no unmortgage. Jail cards stay.\par}
\Needspace{8\baselineskip}\subsection*{G2 d216, Sol}
{\small\textit{Applied private note. Full first resolved field.}\par Both houses yield \$150, leaving \$283 and \$167 short. Clearing buildings unlocks yellow mortgages worth \$430; mortgage Marvin plus one other to survive with two yellows mortgaged. Bankruptcy is avoidable.\par}
\Needspace{8\baselineskip}\subsection*{G4 d270, Gemini}
{\small\textit{Applied private note. Full first resolved field.}\par I owe \$750 to p1 with a \$474 shortfall. I must sell my 4 houses on Yellow (\$300 total) and then mortgage two Yellow properties to cover the remaining \$174 shortfall. First, following even-selling rules, I will sell 1 house on Marvin Gardens.\par}
\Needspace{8\baselineskip}\subsection*{G4 d308, Opus}
{\small\textit{Provider-returned explanation. Full first resolved field.}\par I owe 800 against 45 cash, leaving a 755 shortfall. Selling back my 7 houses at half price gets 525, and mortgaging the three red properties adds 340, totaling 910—enough to cover the debt and leave roughly 110 left, though I'd be stripped of all buildings and have everything mortgaged.

My opponent sits at 1564 with houses everywhere, so my odds are slim, but staying in the game beats certain loss, so I'll liquidate. Selling has to happen incrementally, so I'll start by selling one house off Illinois.\par}
\Needspace{8\baselineskip}\subsection*{G5 d263, Gemini}
{\small\textit{Applied private note. Full first resolved field.}\par I owe \$900 rent on Pacific Avenue and currently have \$74 cash. By selling all 3 houses on Park Place (\$300) and all 3 houses on Boardwalk (\$300), I will reach \$674 cash and unlock the ability to mortgage Park Place (\$175) and Boardwalk (\$200), allowing me to raise \$1,049 total and survive with \$149 cash.\par}
\Needspace{8\baselineskip}\subsection*{G6 d185, Opus}
{\small\textit{Provider-returned explanation. Full first resolved field.}\par With \$750 owed and \$188 cash, I'm short \$562. Selling the houses on St Charles, States, and Virginia Avenues for \$350, plus mortgaging Boardwalk for \$200 and those three properties for another \$220, gives me access to \$770 total—more than enough to cover the shortfall.

That leaves me with \$208 afterward, holding nothing but mortgaged properties, but surviving is still better than going bankrupt now—opponent might bust first and hand me second place. I'll proceed: mortgage Boardwalk first, then work through selling houses one at a time starting with Virginia.\par}
\Needspace{8\baselineskip}\subsection*{G8 d257, Gemini}
{\small\textit{Applied private note. Full first resolved field.}\par I owe \$625 to p3 and have a \$317 shortfall. I must first sell my 1 house on Boardwalk for \$100, which will also unlock Park Place and Boardwalk for mortgaging so I can easily cover the remaining \$217 shortfall.\par}
\Needspace{8\baselineskip}\subsection*{G11 d203, Opus}
{\small\textit{Provider-returned explanation. Full first resolved field.}\par With \$76 cash plus selling the orange houses for \$100 and mortgaging Short Line for \$100, I get \$276 without touching the other properties. Adding the remaining mortgages on St James, Tennessee, and New York (\$90+\$90+\$100) brings the total to \$556, enough to cover the \$500 debt with \$56 left over.

Since paying is possible, I'll avoid bankruptcy and start liquidating now, beginning by selling both houses in a single call since the sell plan handles multiple items at once.\par}
\Needspace{8\baselineskip}\subsection*{G12 d190, Sol}
{\small\textit{Applied private note. Full first resolved field.}\par Three houses sell for \$225; Baltic and both light blues mortgage for \$130. With \$53 cash, maximum liquidity is \$408, below \$500 owed. No legal liquidation can save me.\par}
\Needspace{8\baselineskip}\subsection*{G13 d195, Gemini}
{\small\textit{Applied private note. Full first resolved field.}\par I owe \$750 and have \$168 cash (shortfall \$582). Since my non-Dark-Blue properties only mortgage for \$445 total, I must sell the house on Boardwalk (\$100) and mortgage Park Place (\$175) plus \$335 in non-Dark-Blue properties, keeping Boardwalk and Kentucky Avenue unmortgaged. Selling the house on Boardwalk first.\par}
\Needspace{8\baselineskip}\subsection*{G17 d234, Gemini}
{\small\textit{Applied private note. Full first resolved field.}\par I owe \$550 rent on Tennessee Avenue, and my total liquidatable value is only \$530 (5 Dark Blue houses at \$100 each plus \$30 for mortgaging Baltic Avenue). Since I cannot raise \$550, I must declare bankruptcy.\par}
\Needspace{8\baselineskip}\subsection*{G17 d254, Opus}
{\small\textit{Provider-returned explanation. Full first resolved field.}\par I'm \$422 short of the \$550 I owe, with \$128 cash. Selling back my 5 houses (Kentucky 2, Indiana 1, Illinois 2) at \$75 each gives \$375, and mortgaging those three properties adds \$110+\$110+\$120=\$340, totaling \$843 — more than enough to cover the debt.

But houses must be sold evenly, so I can't dump all 5 at once; I should start by selling one each from Kentucky and Illinois to even things out, bringing in \$150 and leaving a \$272 shortfall before continuing to sell and mortgage further. Staying solvent this way is better than bankruptcy, even if I end up with little cash afterward.\par}
\Needspace{8\baselineskip}\subsection*{G18 d309, Gemini}
{\small\textit{Applied private note. Full first resolved field.}\par I owe \$450 rent on Connecticut Avenue with \$12 cash (\$438 shortfall). Selling all 8 Pink houses raises \$400 (bringing cash to \$412), which will allow me to mortgage St. Charles Place (\$70) on the next step to cover the remaining \$38 shortfall and stay in the game.\par}
\Needspace{8\baselineskip}\subsection*{G18 d313, Grok}
{\small\textit{Applied private note. Full first resolved field.}\par Owe p1 \$550 on Tennessee (3 houses); cash \$1. Must not go bankrupt — p1 would inherit built dark blues and I am eliminated. Selling all 5 houses yields only \$500, so I still need to mortgage PP or BW afterward (\textasciitilde{}\$49 gap). Even-build: BW 3/PP 2, so sell 2 BW now (ends 1/2). Next rounds sell the rest, then mortgage the weaker side if forced. Doubles means I re-roll from Tennessee; most landings are mortgaged/safe except NY (\$650) and Chance. Survival > death.\par}
\Needspace{8\baselineskip}\subsection*{G20 d291, Opus}
{\small\textit{Provider-returned explanation. Full first resolved field.}\par I'm \$725 short even after mortgaging Indiana and Illinois for \$230 and selling all the houses on the pink properties for \$450 — that only totals \$680, still under the \$1000 needed. With Kentucky already mortgaged and the red group empty of houses, I'm checking whether mortgaging the pink lots themselves after clearing the houses could close the gap.\par}
\Needspace{8\baselineskip}\subsection*{G20 d327, Sol}
{\small\textit{Applied private note. Full first resolved field.}\par Full liquidation should yield \$1,414 including cash and all mortgages, enough for \$1,400 rent by \$14. Must sell all development and mortgage everything; survive rather than concede. If hotel sales leave four houses, sell those next evenly.\par}


\section{The full question agenda}
{\small
\begin{longtable}{@{}L{.21\linewidth}L{.39\linewidth}L{.32\linewidth}@{}}
\toprule Question & Supported observation & Unmeasured construct\\\midrule\endhead
1. Strategic reasoning & Prospective budgets, resource dependencies and rival-development objections & Optimal policy, regret, exploitability, hidden cognition\\
2. Opponent modeling & Some checked rival cash and building calculations & Calibrated hidden beliefs and preferences\\
3. Negotiation, cooperation and deception & Complete responses; bluff, omission, correction, appeal and withdrawal & Causal tactic effects, fairness and prevalence\\
4. Planning and memory & Stated local budgets followed by dependent payments & Lookahead depth, memory benefit, general long-plan success\\
5. Adaptation and learning & Revised offers and a within-response correction & Durable learning or unfamiliar-opponent transfer\\
6. Economics and risk & Complete debt requirements and actual financing & Optimal reserves, asset values and liquidation quality\\
7. Collective dynamics & Financial-denial appeal and ownership transfer & Stable coalition prevalence and effectiveness\\
8. Evaluation validity & Channels, delivered menus, response conditions and named replay modes & Repaired-interface or alternative-roster performance\\\midrule
A. Talk versus action & Local matches, factual errors and conflicting calls & General explanation/action predictive relation\\
B. General versus multi-agent ability & No independent capability axis & Capability curves or cross-task relation\\
C. Skill, chance and weak opponents & Five seed blocks and cyclic seats & Causal shares or counterfactual winners\\
\bottomrule
\end{longtable}}

\section{Artifacts and replay}
Repository: \url{https://github.com/KushagraBharti/MonopolyBench}. Source commits are \path{42f85183fa772b2ab3759cc49987e5ed6cc9efcc} for G1--G4 and \path{d6808534a831862682a6940f394c6c87bb2691cb} for G5--G20. The five seeds, in block order, are 751774554, 2506298840, 2573692901, 3648314310 and 2766649504. The frozen manuscript packages retain full game/run maps, original-field manifests and saved-record analysis entry points. The review insert's arithmetic is reproduced by \path{check_saved_records.py}; \path{prepare_insert.py} prepares its source-checked quotation and table includes. These scripts import no engine or model client. This local insert is not itself a published checkpoint of the research package.

All twenty saved \path{artifact_replay_report.json} files pass full-event equality, covering 50,918 events with no event exclusions or field removal. Sorted-key JSON changes representation only; a synthetic clock starts at zero and advances 250 ms per event. All twenty \path{state_replay_report.json} files pass comparison of 24,506 state-relevant events after excluding four LLM observation event types and normalizing global identifiers, sequence and time. These are separate comparison objects. We read the reports and do not execute replay. Equality establishes reproducibility under the frozen implementation, not agreement with an independent rules specification.

\begin{thebibliography}{9}
\bibitem{turpin} M. Turpin, J. Michael, E. Perez and S. R. Bowman. Language Models Don't Always Say What They Think: Unfaithful Explanations in Chain-of-Thought Prompting. NeurIPS, 2023.
\bibitem{lanham} T. Lanham et al. Measuring Faithfulness in Chain-of-Thought Reasoning. arXiv:2307.13702, 2023.
\bibitem{guhe} M. Guhe and A. Lascarides. The Effectiveness of Persuasion in The Settlers of Catan. IEEE CIG, 2014.
\bibitem{jehiel} P. Jehiel, B. Moldovanu and E. Stacchetti. How (Not) to Sell Nuclear Weapons. American Economic Review, 86(4):814--829, 1996.
\end{thebibliography}
\end{document}
